% ----------------------------------------------------
% Projeto: UnDF-TeX
% Arquivo: 04-dados.tex
%
% Exemplo de seção opcional para caracterização
% extensa dos dados ou do corpus da pesquisa.
% ----------------------------------------------------

\chapter{Dados e corpus documental}

Este exemplo não realiza pesquisa com participantes. Seu corpus demonstrativo é formado pelos elementos de um trabalho acadêmico, pelas normas técnicas indicadas na documentação do projeto e pelos recursos implementados no template UnDF-TeX.

Em uma pesquisa empírica, esta seção pode caracterizar a amostra ou o conjunto de dados depois de explicar, na metodologia, como a coleta foi planejada. Em uma pesquisa documental, pode descrever a composição do corpus, sua origem, seu período e sua organização.

\section{Composição do corpus}

Para fins de demonstração, o material foi organizado em três grupos: elementos pré-textuais, recursos do desenvolvimento e elementos pós-textuais. Essa divisão permite verificar a presença, a sequência e a apresentação dos componentes do TCC sem atribuir ao exemplo resultados de uma pesquisa inexistente.

\section{Preparação para análise}

Os arquivos foram reunidos por função, compilados e conferidos no documento final. Em um estudo real, esta etapa também deve registrar exclusões, correções, dados ausentes, transformações e demais decisões que alterem o conjunto bruto antes da análise.
